\section{Integral realizations and the limits of prime tests}
\label{rev:arithmetic-details}

\subsection{The finite Clausen theorem}
\label{subsec:finite-clausen-precision}
Write $s=1/d$, with $d\in\{2,3,4,6\}$, and let $U_p,T_p$ be
the canonical truncations of Section~\ref{sec:weighted-cartier-criterion}.
Set
\[
 k_2=k_6=-1,\qquad k_3=-3,\qquad k_4=-2,\qquad
                  \chi_{d,p}=\left(\frac{k_d}{p}\right).
\]

\begin{theorem}[Finite Clausen and reflection]
\label{thm:brstt-polynomial-packet}\label{cor:weighted-clausen-p2}
For every $p>3$, the following are polynomial congruences:
\begin{align}
 T_p(4x(1-x))&\equiv U_p(x)^2\pmod{p^2\Z_p[x]},
                                      \label{eq:finite-clausen-modp2}\\
 U_p(x)&\equiv\chi_{d,p}U_p(1-x)\pmod{p^2\Z_p[x]}.
                                      \label{eq:finite-reflection-p2}
\end{align}
For $A,B\in\Z_p$, their differentiated form is
\[
 (1-2x)W_p(4x(1-x))\equiv
 U_p(x)\bigl(A(1-2x)U_p(x)+2Bx(1-x)U_p'(x)\bigr)\pmod{p^2}.
\]
In particular the weighted congruence requires no inversion until
specialization away from $1-2x=0$.
\end{theorem}
\begin{proof}
The first identity is \cite[\S7, equation (7.4)]{BRSTT}, where it
is attributed to Chisholm--Deines--Long--Nebe--Swisher, Lemma 18.
We explain the reflection step because its polynomial content is
needed over residue-field extensions.

For the elliptic families of \cite[Proposition 7.1]{BRSTT}, that
proposition gives, over every finite field $\mathbf F_q$ of
characteristic greater than three,
\[
       a_q(\widetilde E_d(x))
         =\chi_q(k_d)a_q(\widetilde E_d(1-x))
\]
at smooth parameters.  It does so by a geometric isogeny of degree
$1,3,2,1$ for $d=2,3,4,6$, respectively, followed by a quadratic
twist.  The source families have
\[
\begin{array}{c|c|c}
 d&j(\widetilde E_d(t))&t\text{ corresponding to }E_{1/d,x}\\\hline
 2&256(1-t+t^2)^3/[t^2(1-t)^2]&x\\
 3&27(9-8t)^3/[t^3(1-t)]&1-x\\
 4&64(4-3t)^3/[t^2(1-t)]&1-x\\
 6&432/[t(1-t)]&x.
\end{array}
\]
Their Hasse polynomials are $U_p$, up to nonzero differential
normalizations.  The trace identity proves the reflection
congruence modulo $p$ at all smooth $\mathbf F_p$-parameters.
The terminating Chu--Vandermonde identity gives
$U_p(0)=1$ and $U_p(1)\equiv\chi_{d,p}\pmod p$, including the
two omitted parameters.  Since the reflected difference has
degree less than $p$, this proves its vanishing as a polynomial.

Write $U_p(x)=\chi_{d,p}U_p(1-x)+pN(x)$.  Invariance of the
left side of \eqref{eq:finite-clausen-modp2} under $x\mapsto1-x$
gives
\[
                      2\chi_{d,p}N(x)U_p(1-x)=0
                              \quad\text{in }\mathbf F_p[x].
\]
This integral domain has $U_p\ne0$, so $N=0\bmod p$.
This is the modulus-$p^2$ argument of \cite[Lemma 7.2]{BRSTT}.
Differentiation preserves the ideal $p^2\Z_p[x]$ and yields
the weighted identity by the chain rule.
\end{proof}
The reflected product
$\chi_{d,p}U_p(u)U_p(v)$, with $u+v=1$ and $uv=z/4$, is
symmetric in $u,v$ and therefore belongs to $\Z_p[z]$.
It equals $T_p(z)$ modulo $p^2$.  This descent explains why all
these tests remain valid when the two parameters lie in an
unramified quadratic extension rather than in $\Q_p$.

\subsection{Cartier and Gauss--Manin differentiation}
\label{rev:cartier-calculation}
We give the normalization behind
Lemma~\ref{lem:cartier-hypergeometric-bridge}.  Write
\[
 P(X)=4X^3-g_2(x)X-g_3(x),\qquad m=(p-1)/2,
 \qquad H=[X^{p-1}]P^m,\quad J=[X^{p-2}]P^m.
\]
For the standard basis $\Omega=dX/y$, $\Xi=X\,dX/y$, write
\[
 \nabla_x\Omega=a\Omega+b\Xi,\qquad
 \nabla_x\Xi=c\Omega-a\Xi.
\]
The exact differential reductions defining this connection imply
\begin{equation}\label{eq:hasse-gm-derivative}
                          H'=aH+bJ,\qquad J'=cH-aJ.
\end{equation}
Indeed, multiply either reduction identity by $P^{m-1}$.
Because $m=-1/2$ in characteristic $p$, its exact correction
becomes $\partial_X(RP^m)$, whose coefficient of $X^{p-1}$ is
zero.  Thus the Gauss--Manin system survives coefficient
extraction.  In particular $H$ satisfies
\[
                 x(1-x)H''+(1-2x)H'-s(1-s)H=0.
\]

For $s=1/6,1/4,1/3$, the model coefficients give $\deg_xH<p$.
At $x=0$ their common cubic is $(X-1)(2X+1)^2$.  The coefficient
of $X^{p-1}$ in a polynomial of degree less than $2p-1$ is
translation invariant modulo $p$: the possibly interfering
coefficients have $\binom{j}{p-1}=0\bmod p$ for
$p\le j<2p-1$.  Translation by $-1/2$ therefore gives
$H(0)=(-6)^m$.  The Gauss equation's recurrence
\[
       (n+1)^2h_{n+1}=(n+s)(n+1-s)h_n\quad(0\le n<p-1)
\]
then identifies $H=(-6)^mU_p$.
For $s=1/2$, the affine change
$X=3T/2-(1+x)/2$ turns the cubic into
$(27/2)T(T-1)(T-x)$.  The Legendre Hasse polynomial
$\sum_{n=0}^m\binom mn^2x^n$ and the same coefficient
invariance give the identical normalization.  Hence, in all four
families,
\begin{equation}\label{eq:hasse-hypergeometric}
                  H=(-6)^mU_p.
\end{equation}

Cartier coefficient extraction now gives
\[
 \operatorname{Cartier}(\nu_{A,B})
 =\bigl((A+BQa)H+BQbJ\bigr)^{1/p}\Omega
 =\bigl((-6)^m(AU_p+BQU_p')\bigr)^{1/p}\Omega.
\]
There is only one coefficient because
$\deg_X((A+BQa+BQbX)P^m)<2p-1$; the notation incorporates
the usual Frobenius twist.  Together with the conjugate-filtration
sequence, this proves the full Cartier statement over arbitrary
perfect residue fields.  It also shows why a change of period
normalization cannot be ignored in the finite-polynomial test.

\subsection{Projective Frobenius and rational-prime density}
\label{rev:qcurve-proof}
We supply the Galois details used in
Theorem~\ref{thm:B-qcurve-density}.

\begin{lemma}[Supersingular projective Frobenius]
\label{lem:B-supersingular-frobenius}
For a supersingular elliptic curve over $\mathbf F_{p^f}$,
$p\ge5$, Frobenius on any auxiliary Tate module has projective
order dividing six.
\end{lemma}
\begin{proof}
Write the Frobenius polynomial as $T^2-aT+p^f$.
Supersingularity gives slopes $1/2,1/2$, so
$v_p(a)\ge f/2$; the Hasse--Weil bound gives
$|a|\le2p^{f/2}$.  If $f$ is odd these inequalities force $a=0$,
and the eigenvalue ratio is $-1$.  If $f$ is even, then
$a/p^{f/2}\in\{-2,-1,0,1,2\}$.  The middle three values give
eigenvalue ratios of order dividing six.  At either endpoint
Frobenius is scalar: its characteristic equation would otherwise
give a nonzero nilpotent in the rational endomorphism division
algebra of an elliptic curve.  This proves the assertion.
\end{proof}

For the non-CM $\Q$-curve, choose a finite Galois field of
definition $K/\Q$, write $d=[K:\Q]$, and choose isogenies
$\mu_\sigma:{}^\sigma E\to E$, with $\mu_1=1$.
They need not be defined over $K$.  For every sufficiently large
auxiliary prime $\lambda$, the maps
\[
 A_g=\mu_{g|_K}\circ g:E[\lambda]\longrightarrow E[\lambda]
\]
are invertible.  The rational Hom spaces between conjugate
non-CM curves are one-dimensional, so
$A_gA_h=c(g,h)A_{gh}$ with $c(g,h)\in\Q^\times$.
Its square is a quotient of the finitely many isogeny degrees;
after excluding their prime divisors it is a $\lambda$-adic
unit.  Projectivization therefore gives a continuous representation
\[
 \overline\rho_\lambda:G_\Q\longrightarrow
                          \operatorname{PGL}_2(\mathbf F_\lambda).
\]
Continuity follows on adjoining the isogeny coefficients and
torsion coordinates to a finite Galois extension.  On $G_K$ this
is the usual projective representation.  Serre's open image
theorem \cite{Serre}, in its number-field form, says that the
residual image on $G_K$ contains
$\operatorname{SL}_2(\mathbf F_\lambda)$ for every sufficiently
large $\lambda$.  Thus the projective image contains
$\operatorname{PSL}_2(\mathbf F_\lambda)$.

Fix such a $\lambda$ and omit its finite ramification and
bad-reduction sets.  At a supersingular rational prime, let $g$
be Frobenius and let $f\mid d$ be the residue degree in $K$.
Then $g^f\in G_K$ acts as Frobenius over $\mathbf F_{p^f}$.
The lemma gives
\[
                 \overline\rho_\lambda(g)^{6d}=1.
\]
Put $M=6d$ and take $\lambda>M$.  An element killed by $M$ is
semisimple.  The split and nonsplit torus types have orders
$\lambda-1$ and $\lambda+1$, and each contains at most $M$
such elements.  Hence at most $2M$ conjugacy classes occur,
each of size at most $\lambda(\lambda+1)$.  Since
$|\operatorname{PSL}_2(\mathbf F_\lambda)|
=\lambda(\lambda^2-1)/2$, their proportion in the residual
image is at most
\[
                         \frac{4M}{\lambda-1}.
\]
Chebotarev identifies the density of each fixed finite-image
condition with its cardinality proportion.  The upper natural
density of our fixed supersingular prime set is bounded by this
quantity for every sufficiently large $\lambda$, hence is zero.
No uniform Chebotarev estimate is used: the prime-counting limit
is taken before $\lambda$ tends to infinity.

\subsection{The analytic order in Tang's inequality}
\label{subsec:B-growth-sharp}
For completeness the order convention in
Theorem~\ref{thm:B-Tang-external} is
\[
 T_\gamma(r)=\int_0^r\frac{dt}{t}
       \int_{\|z\|<t}\gamma^*\eta\wedge\omega^{d-1},
 \qquad\omega=dd^c\log\|z\|^2,
 \qquad\rho_\gamma=\limsup_{r\to\infty}
                         \frac{\log T_\gamma(r)}{\log r},
\]
where $\eta$ is the curvature of an ample Hermitian line bundle
on a projective compactification.  The order is independent of
these choices in the setting used here.

For an elliptic universal extension the order-two estimate has
a direct geometric explanation.  Write
$E(\mathbb C)=\mathbb C/\Lambda$, and use its Weierstrass
functions $\wp,\zeta,\sigma$ and quasiperiod homomorphism
$\eta_\lambda$.  After rescaling the additive coordinate,
\[
 G(\mathbb C)=\mathbb C^2/
           \{(\lambda,\eta_\lambda):\lambda\in\Lambda\}.
\]
A complementary line has exponential map
$t\mapsto[(at,bt)]$ with $a\ne0$.  On the torsor above
$E\setminus\{0\}$ the rational coordinates
$\wp(u),\wp'(u),v-\zeta(u)$ are invariant under the lattice;
the addition theorem for $\zeta$ gives their rational group law.
Since $\sigma$ is entire of order two, its logarithmic
derivatives, and hence these three pulled-back coordinates, have
Nevanlinna order at most two.  Rational characteristic inequalities
give the same bound in any projective embedding.

Conversely the projection $t\mapsto at\bmod\Lambda$ pulls an
ample invariant curvature form on $E$ back to a positive constant
multiple of $|a|^2i\,dt\wedge d\bar t$.  Its characteristic is
therefore of order exactly two.  Comparison of ample bundles on
a projective-bundle compactification of $G$ bounds this
characteristic by that of the original map.  Thus the order is
exactly two.  The threshold $3/4$ cannot be improved solely by
sharpening this growth estimate; this is not a claim that the
arithmetic threshold itself is optimal.

\subsection{Unweighted zeros and supersingularity}
\label{rev:unweighted-tests}
\begin{proposition}[An arithmetic classification of parameters]
\label{prop:levels123-density}\label{prop:density}
\label{prop:infinite-congruences}
For $|C|\ge R_\ell$, $C\ne R_\ell$, and all sufficiently good
primes $p$, one has
\[
 T_p(R_\ell/C)=0\bmod p
 \quad\Longleftrightarrow\quad T_p(R_\ell/C)=0\bmod p^2
 \quad\Longleftrightarrow\quad E_{s,x}\text{ is supersingular}.
\]
Their rational-prime set has natural density $1/2$ for exactly
the 36 level--denominator pairs in the standard tables, and
density zero otherwise.  At each admitted prime $W_p=0\bmod p$
for every coefficient pair.  At level four the set is infinite
for every parameter in this domain, including the non-CM ones.
\end{proposition}
\begin{proof}
Work in the unramified quadratic algebra containing the roots
of $x^2-x+(R_\ell/C)/4$.  The Hasse criterion and the finite
Clausen square show that vanishing modulo $p$ is equivalent to
supersingularity, and that it implies divisibility by $p^2$.
Divisibility descends to $\Z_p$ by unramifiedness; the reverse
implication is immediate.  Differentiating the square gives
$W_p=0\bmod p$ when its Hasse factor vanishes.  The density
assertion now follows from Theorem~\ref{thm:B-qcurve-density}
and the finite CM classification.

For infinitude at level four, descend geometrically to the
rational curve
\begin{equation}\label{lem:rational-descent}
 \mathcal E_z:\quad y^2=4X^3-3(1-z/4)(1-z)X
                                      -(8+z)(1-z)^2/8.
\end{equation}
Its discriminant is $729z^2(1-z)^3/64$ and
$j=64(4-z)^3/z^2$, which agrees with that of the Legendre fiber
when $z=4x(1-x)$.  Hence geometric supersingularity agrees
outside finitely many primes.  Elkies' theorem \cite{Elkies}
states that every elliptic curve over $\Q$ has infinitely many
supersingular primes, proving the claim.
\end{proof}

For example, $C=72$ at level four has $j=21952/9$ and is non-CM,
yet has infinitely many such zeros.  Thus infinitude alone cannot
replace positive density in the unweighted criterion.
Nor can finitely many tests distinguish the CM rows: fix finitely
many $p\ge5$ and exponents $e_p$, and put $M=\prod p^{e_p}$.
Every $C=256+kM>4096$ has, with $(A,B)=(1,6)$, the same
weighted truncations of every length modulo each $p^{e_p}$ as
$C=256$, term by term.  Nevertheless
\[
                  j(C)=C^2-48C+768-4096/C
\]
is nonintegral, so all these parameters are non-CM.

\subsection{Higher precision and the missing comparison}
\label{subsec:precision-versus-transfer}
At an ordinary prime, the weighted factorization defines a
canonical local slope, but its lift modulo $p^2$ differs from
the rational CM coefficient slope.

\begin{proposition}[The local slope]
\label{prop:local-slope-p2}
Let $p>3$ and $z\in\Z_p$. If $z,1-z$ are units and the fiber is ordinary, then
\[
                  r_p=-Q(x)\frac{U_p'(x)}{U_p(x)}\pmod{p^2}
\]
is independent of the root $x$ and lies in $\Z/p^2\Z$.
For $A,B\in\Z_p$ with $B$ a unit,
$W_p(z)=0\bmod p^2$ if and only if $A/B=r_p\bmod p^2$.
Modulo $p$ this is the complementary Cartier-kernel line.
\end{proposition}
\begin{proof}
The algebra $\mathcal O=\Z_p[x]/(x^2-x+z/4)$ is finite etale;
its involution has fixed ring $\Z_p$.  Ordinarity makes $U_p(x)$
a unit.  Reflection and its derivative give
$U_p'(x)\equiv-\chi_{d,p}U_p'(1-x)\bmod p^2$.
Since $Q(1-x)=-Q(x)$, the quotient defining $r_p$ is invariant,
hence rational modulo $p^2$; invariants commute with reduction
because two is invertible.  The factorization becomes
$W_p\equiv BU_p^2(A/B-r_p)\bmod p^2$, proving the assertions.
\end{proof}

By contrast, \cite[Theorem 2.1]{BRSTT}, restating the theorem of
Chisholm--Deines--Long--Nebe--Swisher, assumes both CM of the
source fiber and its identity
$(1+\beta\thetaop)F_{1/d}(z)=\delta/\pi$ with $|z|<1$.
Under its good-reduction, unramifiedness, and $p$-adic-unit
hypotheses, its conclusion is
\begin{equation}\label{eq:cm-weighted-first-order}
 T_p(z)+\beta zT_p'(z)
       \equiv\epsilon_p\left(\frac{1-z}{p}\right)p\pmod{p^2},
 \qquad
 \epsilon_p=\begin{cases}1&\text{ordinary},\\-1&\text{supersingular}.
                         \end{cases}
\end{equation}
We use this source only in its stated strict convergence range.
For a CM identity with $AB\ne0$, putting $\beta=B/A$ therefore
gives, outside finitely many primes,
\begin{equation}\label{eq:cm-Wp-nonzero-p2}
 W_p(z)\equiv A\epsilon_p\left(\frac{1-z}{p}\right)p\pmod{p^2}.
\end{equation}
Thus $v_p(W_p)=1$.  At ordinary primes the exact first-order drift
of the local slope is
\[
 r_p\equiv\frac AB\left(1-
       \frac{\epsilon_p\left(\frac{1-z}{p}\right)p}{U_p(x)^2}
                   \right)\pmod{p^2}.
\]
For the valid row $(A,B,C)=(1,6,256)$, one obtains
$W_7(1/4)=42\bmod49$, $r_7=20\bmod49$, whereas
$1/6=41\bmod49$.  A transfer demanding $W_p=0\bmod p^2$
would therefore fail even for a tabulated identity.  Likewise the
universal finite Clausen congruence cannot be strengthened to
modulus $p^3$: at $s=1/2$ and $p=5$ its discrepancy is
$75x^6\bmod125$.

\begin{lemma}[The complex series does not sum directly $p$-adically]
\label{lem:padic-divergence}
At every prime $p\nmid C$, a pair-primitive standard series
diverges as a $p$-adic series.
\end{lemma}
\begin{proof}
Choose $e\in\{0,1\}$ with $A+Be\ne0\bmod p$.  For
$n=e+p^r$ and $p^r>6$, the base-$p$ digit blocks of $kn$ are
disjoint for $1\le k\le6$.  Legendre's formula for factorial
valuations, applied to the balanced factorial ratio $u_\ell$,
therefore gives
\[
 v_p(u_\ell(e+p^r))=v_p(u_\ell(e))+v_p(u_\ell(1)).
\]
The linear weight and $C^{-n}$ are units.  Along this subsequence
the summand valuation is constant and finite, so the terms do not
tend to zero.
\end{proof}

Raw finite truncations also need not equal Frobenius unit roots
modulo $p^2$.  At $s=1/2$, $x=1/3$, $z=8/9$, the fiber is
non-CM and has trace $-2$ at $5$.  The unit root of
$T^2+2T+5$ is $13\bmod25$, while
\[
                  U_5(1/3)=8,\qquad T_5(8/9)=14
                                      \quad\bmod25,
\]
and $13^2=19\bmod25$.  The unit-root statement in
\cite[Theorem 2.2]{BRSTT} retains its CM premise.  In that work
the CM action first selects a second-kind eigenvector, after which
a compatible Frobenius lift yields the supercongruence
\cite[\S2.1, equation (2.3), and Lemma 2.3]{BRSTT}.
Neither the universal polynomial identity nor this CM construction
provides the eigenvector from an isolated complex period product.

\subsection{Exact compatibility, finite precision, and a common tensor}
\begin{proposition}[Finite precision does not force CM]
\label{prop:one-prime-precision-c}
There are non-CM elliptic curves over $\Q$ with a common good
ordinary reduction at $5$ whose Hodge-preservation obstruction
in Theorem~\ref{thm:one-prime-c} is arbitrarily small $5$-adically
in fixed analytic local frames.
\end{proposition}
\begin{proof}
Consider $E_N:y^2=x^3-x+5^N$, $N\ge2$.  Its discriminant is
$16(4-27\cdot5^{2N})$, a $5$-adic unit, and its common reduction
has eight $\mathbf F_5$-points, hence ordinary trace $-2$.
Its invariant
$j(E_N)=6912/(4-27\cdot5^{2N})$ is a nonzero rational number
of absolute value less than one, hence nonintegral and non-CM.
The curves tend $5$-adically to the CM curve $y^2=x^3-x$.
On a sufficiently small deformation disk, the family
crystalline--de Rham comparison identifies the fixed special-fiber
isocrystal with the analytic de Rham bundle and makes the Hodge
line analytic \cite[\S4.6]{MaulikPoonen}.  The obstruction
$F^1H\to H/F^1H$ therefore tends to zero.  It is nonzero at
every $E_N$ by Theorem~\ref{thm:one-prime-c}.
\end{proof}

The distinction between local and common invariants is equally
sharp on the Galois side.
\begin{proposition}[A common Galois tensor]\label{prop:galois-tensor}
For an elliptic curve $E/K$ and an auxiliary prime $\lambda$,
$E$ has geometric CM if and only if
\[
           \operatorname{Sym}^2V_\lambda(E)\otimes
                       \det(V_\lambda(E))^{-1}
\]
has a nonzero invariant over some finite extension of $K$.
\end{proposition}
\begin{proof}
For every two-dimensional vector space, polarization gives
\[
             \operatorname{Sym}^2V\otimes\det(V)^{-1}
                           \simeq\operatorname{End}^0(V),
\]
where the superscript denotes trace zero.  Faltings' endomorphism
theorem \cite{Faltings} identifies the invariant endomorphisms
over a finite extension $L/K$ with
$\operatorname{End}_L(E)\otimes\Q_\lambda$.
Its trace-zero part is zero for a non-CM curve and nonzero once
the CM endomorphisms are defined.  This proves both directions.
\end{proof}
For every elliptic curve over $\Q$, however, each good Frobenius
on the untwisted symmetric square has eigenvalues
$r_p^2,r_ps_p,s_p^2$, with $r_ps_p=p$.  The Tate eigenvalue is
therefore present at every prime, including for non-CM curves;
its eigenline need not be common to the different Frobenii.
Factoring off these local eigenvalues cannot by itself prove an
analytic pole of an $L$-function, since the remaining factor may
have a zero.  A complex scalar relation, a collection of local
eigenvalues, and a common Galois tensor are distinct data.
